\documentclass[11pt]{article}
\usepackage[margin=1in]{geometry}
\usepackage{amsmath,amssymb,amsthm}
\usepackage{xurl}
\usepackage{hyperref}
\sloppy
\let\oldus\_ \renewcommand{\_}{\oldus\allowbreak}
\newtheorem{theorem}{Theorem}
\newtheorem{lemma}[theorem]{Lemma}
\theoremstyle{definition}
\newtheorem{definition}[theorem]{Definition}
\newcommand{\fl}[1]{\left\lfloor #1 \right\rfloor}

\title{Proof of conjecture 00000002139:\\ the floor-sum recursion follows the Euclidean algorithm and makes $O(\log m)$ calls}
\author{Jackmeson1}
\date{2026-10-06}

\begin{document}
\maketitle

\section{The conjecture and its reading}

\textbf{Statement (English, verbatim).} Definition: The floor-sum is the summation functional of
division with remainder. Conjecture: The complexity of its recursion is logarithmic; and the
recursion path is the Euclidean algorithm. (floor-sum Euclidean recursion)
(The Chinese text states the same.)

\textbf{The object.} The floor-sum is the function
\[
  F(n,m,a,b)=\sum_{i=0}^{n-1}\fl{\frac{a i+b}{m}},\qquad n\ge 0,\ m\ge 1,
\]
the sum of the quotients of the divisions with remainder of $ai+b$ by $m$. This is the function
\texttt{floor\_sum(n, m, a, b)} of the AtCoder Library (ACL), whose documentation says: ``It returns
$\sum_{i = 0}^{n - 1} \lfloor \frac{a \times i + b}{m} \rfloor$'', with constraints
``$0 \leq n < 2^{32}$'', ``$1 \leq m < 2^{32}$'' and complexity ``$O(\log m)$''
[ACL-doc]. We first take $a,b\ge 0$; signed $a,b$ are handled in Section~\ref{sec:signed}.

\textbf{The recursion (a reading choice).} The text does not spell out ``its recursion''. We take
the standard Euclid-like recursion, as implemented in ACL's \texttt{floor\_sum\_unsigned}
[ACL-code]: on modulus $m$ and slope $a$, add $\binom{n}{2}\fl{a/m}+n\fl{b/m}$, replace $a,b$ by
$a\bmod m$, $b\bmod m$, put $y=(a\bmod m)n+(b\bmod m)$, and continue with
$(n,m,a,b)\leftarrow(\fl{y/m},\ a\bmod m,\ m,\ y\bmod m)$ (``\texttt{std::swap(m, a)}''). We formalize
two stopping rules. For both we prove that they compute the floor sum and make at most $2\log_2 m+1$ calls; the exact trace equality, the terminal gcd and the Fibonacci lower bound are proved for $\mathtt{fsRec}$ only, while for $\mathtt{aclRec}$ (which may exit early) we prove that its path is a prefix of the Euclid trace:
\begin{itemize}
\item $\mathtt{fsRec}$ stops when the reduced slope $a\bmod m$ is $0$. This is the rule of the
  OI-wiki account of the ``Euclid-like algorithm'', which says that $(a,c)$ are repeatedly reduced
  and swapped ``until $a=0$'' and that this resembles Euclid's algorithm on $(a,c)$ [OI-wiki];
\item $\mathtt{aclRec}$ stops when $y<m$. This is exactly the ACL loop (\texttt{if (y\_max < m) break;}).
\end{itemize}

\textbf{The two claims.} (i) \emph{The recursion path is the Euclidean algorithm}: the list of
(modulus, slope) pairs of the successive calls is the list of argument pairs of Euclid's algorithm
$\gcd(x,y)\to\gcd(y\bmod x,x)\to\cdots$ started at $(m,a)$ and stopped when the first argument is
$0$. Its last modulus is $\gcd(m,a)$. For $\mathtt{fsRec}$ this is an equality of lists. ACL exits
early, so for $\mathtt{aclRec}$ the path is a prefix of the Euclid path.
(ii) \emph{The complexity is logarithmic}: each call does a fixed number of arithmetic operations,
so we measure complexity as the number of calls, in the unit-cost model. We prove that it is at most
$2\log_2 m+1$ and at most $2\log_2\min(a,m)+2$, whatever $n$ and $b$ are. We also prove that it is
at least $\log_2 m+1$ on an infinite family. So the worst case is $\Theta(\log m)$.

\textbf{Conventions.} $\mathbb N\ni 0$. Lean's $/$ and $\%$ on $\mathbb N$ are floor division and
remainder; we prove that the $\mathbb N$-valued $F$ equals the sum of genuine floors
$\lfloor (ai+b)/m\rfloor$ of rationals. $\log_2$ is \texttt{Nat.log 2}, the integer part of the
binary logarithm. For $m=0$ the Lean definitions return $0$ and the empty path; the theorems below
assume $m\ge1$.

\section{Definitions (Lean namespace \texttt{C2139})}

\begin{definition}
$\mathtt{floorSum}\ n\ m\ a\ b=\sum_{i<n}(a i+b)/m$ on $\mathbb N$, and
$\mathtt{floorSumInt}\ n\ m\ a\ b=\sum_{i<n}\lfloor((ai+b:\mathbb Z):\mathbb Q)/m\rfloor$ for
$a,b\in\mathbb Z$.
\end{definition}

\begin{definition}
$\mathtt{euclidTrace}\ x\ y=[\,]$ if $x=0$, and otherwise
$(x,y)::\mathtt{euclidTrace}\ (y\bmod x)\ x$. These are the argument pairs of the calls of Euclid's
algorithm, with the same recursion as Mathlib's \texttt{Nat.gcd\_rec}: $\gcd(x,y)=\gcd(y\bmod x,x)$.
\end{definition}

\begin{definition}
$\mathtt{fsRec}$ and $\mathtt{aclRec}$ return a pair (value, path), where the path lists the
(modulus, slope) pair of every call. With $c=\fl{n(n-1)/2}\cdot\fl{a/m}+n\fl{b/m}$ and
$y=(a\bmod m)n+(b\bmod m)$:
$\mathtt{fsRec}\ n\ m\ a\ b=(c,[(m,a)])$ if $a\bmod m=0$, and otherwise
$(c+v,\ (m,a)::p)$ where $(v,p)=\mathtt{fsRec}\ \fl{y/m}\ (a\bmod m)\ m\ (y\bmod m)$.
$\mathtt{aclRec}$ is the same with the test $y<m$ in place of $a\bmod m=0$.
Both terminate by well-founded recursion on $m$, since $a\bmod m<m$.
\end{definition}

\section{Proof}

\begin{lemma}[\texttt{floorSum\_reduce}]
For $m\ge1$: $F(n,m,a,b)=\fl{n(n-1)/2}\fl{a/m}+n\fl{b/m}+F(n,m,a\bmod m,b\bmod m)$.
\end{lemma}
\begin{proof}
$ai+b=\big((a\bmod m)i+(b\bmod m)\big)+m\big(\fl{a/m}i+\fl{b/m}\big)$, so
$\fl{(ai+b)/m}=\fl{a/m}i+\fl{b/m}+\fl{((a\bmod m)i+b\bmod m)/m}$. Sum over $i<n$, using
$\sum_{i<n}i=\fl{n(n-1)/2}$.
\end{proof}

\begin{lemma}[\texttt{count\_eq}, \texttt{floorSum\_swap}]
If $c\ge1$, then $\#\{k<N: c(k+1)\le x\}=\min(N,\fl{x/c})$. If $a\ge1$ and $b<m$, then with $y=an+b$,
\[ F(n,m,a,b)=F(\fl{y/m},\,a,\,m,\,y\bmod m). \]
\end{lemma}
\begin{proof}
The count follows by induction on $N$, since $c(N+1)\le x\iff N+1\le\fl{x/c}$. For the swap, let
$N=\fl{y/m}$ and $B=y\bmod m$, so $mN+B=an+b$. For $i<n$ we have $\fl{(ai+b)/m}\le N$, so by the count
$F(n,m,a,b)=\sum_{i<n}\sum_{k<N}[m(k+1)\le ai+b]$. For $k<N$ we have
$mk+B\le mN+B-m=an+b-m<an$, so $\fl{(mk+B)/a}\le n$ and
$F(N,a,m,B)=\sum_{k<N}\sum_{i<n}[a(i+1)\le mk+B]$. Exchange the two sums and reflect both indices,
$k\mapsto N-1-k$ and $i\mapsto n-1-i$. The two indicators then agree:
$m(N-k)\le a(n-1-i)+b \iff a(i+1)\le mk+B$, because $m(N-k)+mk=mN$, $a(n-1-i)+a(i+1)=an$ and
$mN+B=an+b$. This counts the lattice points under the line by rows and by columns.
\end{proof}

\begin{lemma}[\texttt{euclidTrace\_last}, \texttt{euclidTrace\_length}, \texttt{euclidTrace\_length\_min}]
Let $x\ge1$. The last pair of $\mathtt{euclidTrace}\ x\ y$ has first entry $\gcd(x,y)$. The trace has
length at most $2\log_2 x+1$, and at most $2\log_2\min(y,x)+2$.
\end{lemma}
\begin{proof}
For the gcd, use strong induction and $\gcd(x,y)=\gcd(y\bmod x,x)$. If $y\bmod x=0$, then
$\gcd(x,y)=x$. For the first bound, let $r=y\bmod x$. If $r=0$, the length is $1$. Otherwise the
trace is $(x,y),(r,x)$ followed by $\mathtt{euclidTrace}\ s\ r$, where $s=x\bmod r$. Since $r<x$
gives $\fl{x/r}\ge1$, we have $s+r\le x$ and $s<r$, so $2s<x$. If $s=0$, the length is
$2\le 2\log_2x+1$, because $x\ge2$. Otherwise $\log_2 s+1=\log_2(2s)\le\log_2x$, and induction gives
length $\le 2+2\log_2 s+1\le2\log_2x+1$. For the second bound, apply the first one to the tail
$\mathtt{euclidTrace}\ r\ x$, using $r\le\min(y,x)$.
\end{proof}

\begin{theorem}[\texttt{fsRec\_spec}, \texttt{aclRec\_spec}]
For all $n,a,b$ and all $m$: $\mathtt{fsRec}\ n\ m\ a\ b=(F(n,m,a,b),\ \mathtt{euclidTrace}\ m\ a)$.
Also $\mathtt{aclRec}$ has value $F(n,m,a,b)$, and its path is a prefix of
$\mathtt{euclidTrace}\ m\ a$.
\end{theorem}
\begin{proof}
Use strong induction on $m$. The case $m=0$ is trivial. For $m\ge1$, apply the reduction lemma. If
$a\bmod m=0$, the remaining sum $F(n,m,0,b\bmod m)$ is $0$, and
$\mathtt{euclidTrace}\ m\ a=[(m,a)]$. Otherwise apply the swap with slope $a\bmod m\ge1$ and offset
$b\bmod m<m$, then the induction hypothesis at the modulus $a\bmod m<m$. The tail path is
$\mathtt{euclidTrace}\ (a\bmod m)\ m$, and $\mathtt{euclidTrace}\ m\ a=(m,a)::$ that tail. For
$\mathtt{aclRec}$, if $y<m$, every term $\fl{((a\bmod m)i+b\bmod m)/m}$ with $i<n$ is $0$, and
$[(m,a)]$ is a prefix. If $y\ge m$, then $a\bmod m\neq0$, and the argument is as before.
\end{proof}

\begin{theorem}[\texttt{euclidTrace\_fib}, \texttt{fsRec\_fib}]
For every $k$, $n$ and $b$, $\mathtt{fsRec}\ n\ F_{k+2}\ F_{k+3}\ b$ makes exactly $k+1$ calls, and
$\log_2F_{k+2}+1\le k+1$. Here $F_j$ is the Fibonacci number \texttt{Nat.fib j}.
\end{theorem}
\begin{proof}
We have $F_{k+4}\bmod F_{k+3}=F_{k+2}$, because $F_{k+2}<F_{k+3}$. Hence the Euclid trace from
$(F_{k+3},F_{k+4})$ is one pair longer than the trace from $(F_{k+2},F_{k+3})$. Also $F_{k+2}\le2^k$.
\end{proof}

\section{Signed $a,b$}\label{sec:signed}
ACL's \texttt{floor\_sum} for signed $a,b$ first replaces them by their residues mod $m$ [ACL-code].
We prove (\texttt{floorSumInt\_normalize}):
\[\mathtt{floorSumInt}\ n\ m\ a\ b=\lfloor a/m\rfloor\fl{n(n-1)/2}+n\lfloor b/m\rfloor+F(n,m,a\bmod m,b\bmod m).\]
Theorem \texttt{floorSumInt\_recursion} then states that after this one step $\mathtt{fsRec}$
computes the signed floor-sum, its path is $\mathtt{euclidTrace}\ m\ (a\bmod m)$, and it makes at
most $2\log_2m+1$ calls.

\section{Lean}
The main theorem is \texttt{C2139.floorSum\_recursion\_euclid} (assuming $m\ge1$). It states the
conjunction of: $F$ equals the genuine floor-sum (\texttt{floorSum\_cast}); $\mathtt{fsRec}$ computes
$F$; its path equals $\mathtt{euclidTrace}\ m\ a$; the last modulus is $\gcd(m,a)$; the number of
calls is at most $2\log_2m+1$ and at most $2\log_2\min(a,m)+2$; $\mathtt{aclRec}$ computes $F$, and
its path is a prefix of $\mathtt{euclidTrace}\ m\ a$ of length at most $2\log_2m+1$. Companion
theorems: \texttt{C2139.floorSumInt\_recursion} (signed) and \texttt{C2139.fsRec\_fib} (tightness).
The file \texttt{lean/Conjecture2139/Basic.lean} imports only Mathlib. It has no \texttt{sorry}, no
\texttt{native\_decide} and no added axioms. \texttt{\#print axioms} reports only
\texttt{propext}, \texttt{Classical.choice} and \texttt{Quot.sound} for all three theorems.

\section*{Sources}
\begin{itemize}
\item[{[ACL-doc]}] AtCoder Library, math documentation, section floor\_sum.
\url{https://github.com/atcoder/ac-library/blob/master/document_en/math.md}
\item[{[ACL-code]}] AtCoder Library source, \texttt{floor\_sum\_unsigned} in \texttt{atcoder/internal\_math.hpp} and \texttt{floor\_sum} in \texttt{atcoder/math.hpp}.
\url{https://github.com/atcoder/ac-library/blob/master/atcoder/internal_math.hpp}
\item[{[OI-wiki]}] OI-wiki, ``Euclid-like algorithm'' (in Chinese).
\url{https://github.com/OI-wiki/OI-wiki/blob/master/docs/math/number-theory/euclidean.md}
\end{itemize}

\end{document}
